\chapter{Introduction et Contexte}

\section{Contexte du projet}

Le secteur hospitalier fait face à une transformation numérique profonde. La gestion
manuelle des dossiers patients, la prise de rendez-vous par téléphone et la coordination
entre les différents services génèrent des inefficacités coûteuses et des risques d'erreurs
médicales. \textbf{MedSys} répond à ce besoin en proposant une plateforme numérique
intégrée, pensée pour les établissements de santé de taille moyenne à grande.

\section{Objectifs}

\begin{itemize}[itemsep=4pt]
  \item \textbf{Centraliser} les dossiers médicaux électroniques dans un système sécurisé
    et accessible aux professionnels de santé autorisés.
  \item \textbf{Simplifier} la prise de rendez-vous pour les patients, avec ou sans compte,
    grâce à une interface web intuitive.
  \item \textbf{Automatiser} les notifications (confirmations, rappels J-1, libération
    de créneaux) par email et en temps réel.
  \item \textbf{Outiller} la direction avec des tableaux de bord analytiques complets.
  \item \textbf{Gérer} les ressources humaines du personnel hospitalier (plannings,
    absences, équipes).
  \item \textbf{Sécuriser} l'ensemble des accès via JWT, RBAC et audit logs.
\end{itemize}

\section{Périmètre fonctionnel}

\begin{table}[H]
\centering
\renewcommand{\arraystretch}{1.4}
\begin{tabularx}{\textwidth}{|l|X|l|}
\hline
\rowcolor{tabhead}\color{white}\textbf{Module} &
\color{white}\textbf{Fonctionnalités principales} &
\color{white}\textbf{Utilisateurs} \\
\hline
ms-auth & Inscription, connexion, JWT, refresh token, RBAC, reset mdp, audit &
Patient, Admin \\
\hline
\rowcolor{roweven}
ms-patient-personnel & Dossier médical, consultations, ordonnances, analyses, PDF,
QR Code, messagerie, stats & Patient, Médecin, Direction \\
\hline
MedicalAppointments & Créneaux, réservation, liste d'attente, notifications, SignalR &
Patient, Médecin, Secrétaire \\
\hline
\rowcolor{roweven}
ms-personnel & Plannings, absences, équipes, disponibilités, demandes &
Personnel, Direction \\
\hline
medsys-web & Interface unifiée, tableaux de bord par rôle, prise de RDV &
Tous les rôles \\
\hline
\end{tabularx}
\caption{Périmètre fonctionnel de MedSys}
\end{table}

\section{Technologies utilisées}

\begin{table}[H]
\centering
\renewcommand{\arraystretch}{1.4}
\begin{tabular}{|l|l|l|}
\hline
\rowcolor{tabhead}\color{white}\textbf{Couche} &
\color{white}\textbf{Technologie} &
\color{white}\textbf{Version} \\
\hline
Backend Java & Spring Boot, Spring Security, Spring Data JPA & 3.x \\
\rowcolor{roweven}
Backend .NET & ASP.NET Core, Entity Framework Core & 8.0 \\
Frontend web & React, Vite, Axios & 18 / 5.x \\
\rowcolor{roweven}
Frontend mobile & React Native (Expo) & SDK 50 \\
Base de données & MySQL, SQL Server, MongoDB & 8 / 2022 / 7 \\
\rowcolor{roweven}
Messagerie async & RabbitMQ & 3.x \\
Temps réel & SignalR (.NET) & 8.0 \\
\rowcolor{roweven}
Auth & JWT (HS256), Refresh Token & -- \\
Conteneurisation & Docker (Dockerfile par service) & 24.x \\
\hline
\end{tabular}
\caption{Stack technologique de MedSys}
\end{table}

\section{Méthodes de modélisation}

Ce rapport présente la modélisation du système selon deux approches complémentaires :

\begin{itemize}[itemsep=3pt]
  \item \textbf{Merise} : modélisation des données à travers trois niveaux —
    MCD (conceptuel), MLD (logique) et MPD (physique). Adapté à la présentation
    de la structure des bases de données relationnelles.
  \item \textbf{UML 2.x} : modélisation comportementale et structurelle — diagrammes
    de cas d'utilisation, de classes, de séquences et de composants.
\end{itemize}
